%%=============================================================================
%% CAPITULO 5 -- CONSIDERACOES FINAIS
%%
%% Texto dummy (exercicio ia-6.1): trechos do pre-projeto do autor (versao 7).
%%=============================================================================

\section{Considerações finais}

No domínio de privacidade, os testes propostos serão entendidos como avaliação
empírica de risco de exposição, e não como demonstração de privacidade absoluta.
A pesquisa buscará identificar sinais de proximidade excessiva, memorização ou
membership disclosure segundo os procedimentos definidos. Garantias formais, como
as proporcionadas por mecanismos de privacidade diferencial, permanecem como
possibilidade de investigação futura. Essa delimitação impede que a ausência de
evidência de exposição nos testes seja interpretada como prova de risco zero.

Outra limitação refere-se à validade externa. A pesquisa será conduzida a partir
de um contexto organizacional específico e de dados associados à gestão de
vulnerabilidades. Ainda que o pipeline seja documentado de forma a permitir
reaplicação, os resultados poderão refletir características particulares do
ambiente estudado, de sua arquitetura tecnológica, de seus processos e de seu
nível de maturidade. Dessa forma, as conclusões serão apresentadas como evidência
de um caso aplicado, e não como demonstração de validade universal para
organizações públicas ou reguladas.
